\section*{Conventions}
\addcontentsline{toc}{section}{Conventions}

\paragraph{Twisted $K$-theory.}
We use nonconnective $K$-theory.
For a central simple algebra $A$ over $F$, we write $K(A)$ for the $K$-theory spectrum of finitely generated projective right $A$-modules.
We write $\per A$ and $\ind A$ for its period and index, that is, the order of $[A]$ in $\Br(F)$ and the degree of the division algebra \(D\) in its class; thus $\ind A=\deg D$.
Morita equivalence identifies $K(A)$ with $K(D)$, where $D$ is the division algebra in the Brauer class of $A$, and all our statements are invariant under this identification.
From now on, we therefore identify $A$ with $D$ in this way and state all results for the division algebra $D$.

The reduced rank \(\Nrd_0\) identifies $K_0(D)$ with the subgroup $\deg D\cdot\Z$ of $K_0(F)=\Z$, the reduced norm $\Nrd_1$ sends the class of $x\in D^\times$ to the determinant of $x\otimes1$ under $D\otimes_FL\simeq M_{\deg D}(L)$ for a splitting field $L$, and $\Nrd_2$ is defined through the $K$-cohomology of the Severi--Brauer variety of $D$ \cite{MerkurjevSuslin1982,Suslin1987}.
For a real place $v$ of $F$, we write the Hilbert symbol $s_v\colon K_2(F)\to\Z/2$ additively: $s_v\{a,b\}=1$ if $a$ and $b$ are both negative at $v$, and $s_v\{a,b\}=0$ otherwise.


\paragraph{Local fields.}
An \emph{$h$-dimensional local field} is a field $F=F_h$ for which there are fields $F_{h-1},\ldots,F_0$ such that $F_0$ is finite and, for $1\le i\le h$, the field $F_i$ is complete discretely valued with residue field $F_{i-1}$.
We call $F_{h-1}$ the \emph{first residue field} and $F_0$ the \emph{last residue field} of $F$, and $(\operatorname{char}F_h,\ldots,\operatorname{char}F_1)$ its \emph{characteristic sequence}.
For a three-dimensional local field $F$, we write $\inv_F\colon H^4(F,\Q/\Z(3))\xrightarrow{\sim}\Q/\Z$ for Kato's isomorphism \cite[Theorem 2]{KatoExistence}.
For a field $F$ of characteristic $p$, we set
\[
 H^{a+1}(F,\Z/p^r(a))=H^1(F,W_r\Omega^a_{\log}),
\]
where $W_r\Omega^a_{\log}$ is the logarithmic de Rham--Witt sheaf.

\paragraph{Spectra.}
We write $\mathrm{Sp}$ for the category of spectra, $\mathbb S$ for the sphere spectrum, $\otimes=\wedge$ for the smash product, and $\Sigma$ for suspension.
The full subcategories $\mathrm{Sp}_{\ge a}$ and $\mathrm{Sp}_{\le b}$ consist of the spectra whose homotopy groups vanish below $a$ and above $b$, respectively, and $\tau_{\ge a}$ and $\tau_{\le b}$ denote the $(a-1)$-connected cover and the $b$th Postnikov truncation.
For an abelian group $M$, we write $HM$ for its Eilenberg--Mac Lane spectrum.
For a spectrum $T$, a prime $\ell$, and an integer $n\ge1$, we write
\(
 T_{\Q}=T\otimes H\Q\), \(T_{(\ell)}=T\otimes\mathbb S_{(\ell)}\), \(T[1/n]\)
for the rationalization, the localization at $\ell$, and the inversion of $n$, and
\(
 T^\wedge_\ell=\operatorname*{holim}_rT/\ell^r
\)
for the derived $\ell$-completion, where $T/\ell^r$ is the cofiber of multiplication by $\ell^r$.
We distinguish $\pi_j(T^\wedge_\ell)$ from the $\ell$-adic completion $\varprojlim_r\pi_j(T)/\ell^r$ of the group $\pi_j(T)$.
The functor $L_{K(1),\ell}$ is Bousfield localization with respect to height-one Morava $K$-theory at $\ell$, and $L_{E(1),\ell}$ is the corresponding localization with respect to height-one Johnson--Wilson theory.
For a finite group $G$ acting on $T$, we write
\[
 T_{hG}=\operatorname*{colim}_{BG}T,\qquad T^{hG}=\operatorname*{lim}_{BG}T
\]
for the homotopy orbits and homotopy fixed points, and for a $G$-module $M$ we write $M_G$ for the ordinary coinvariants.
